\chapter{Viscosupplementation}

\section*{Overview}
Viscosupplementation is the injection of hyaluronic acid, a gel-like lubricant, into a joint to relieve pain and improve function in osteoarthritis. It is most commonly performed in the knee.
\section*{Alternative Names}
Hyaluronic acid injection; hyaluronan injection; knee joint lubricant injection; hylan G-F 20 injection
\section*{Principle}
Injected hyaluronic acid supplements the joint's naturally occurring synovial fluid, restoring its viscoelastic and shock-absorbing properties and reducing friction in the arthritic joint. It may also have anti-inflammatory and cartilage-protective effects.
\section*{History}
The concept of replacing deficient joint fluid developed from research on synovial fluid in the mid-20th century, and hyaluronic acid injections for knee osteoarthritis entered clinical use in the 1980s.
\section*{Why It Is Done}
Ordered for symptomatic knee osteoarthritis in patients with persistent pain who have not responded adequately to conservative measures such as exercise, weight loss, and analgesics, and who wish to delay or avoid surgery.
\section*{Is This a Primary or Secondary Test or Procedure}
A secondary or adjunctive therapeutic procedure, used when first-line osteoarthritis management has been insufficient.
\section*{How It Is Performed}
The joint is identified and the skin is cleaned. Under sterile technique, a needle is inserted into the joint, and one or more injections of hyaluronic acid are given, either as a single dose or as a series separated by a week, depending on the product.
\section*{Preparation}
Informed consent, and the procedure is avoided if the joint shows signs of active infection or significant effusion. Aspiration of joint fluid may be performed first if an effusion is present.
\section*{Contraindications and Precautions}
Contraindications include infection or skin disease over the injection site, known allergy to hyaluronan preparations, and bleeding disorders. Precautions include pregnancy, recent joint injury, and systemic anticoagulation.
\section*{Risks}
Injection-site pain and swelling, transient post-injection joint pain or flare, and rarely infection or bleeding within the joint.
\section*{Aftercare and Recovery}
The patient is advised to rest the joint briefly and avoid strenuous activity for a day or two. Benefit, if achieved, develops over several weeks and may last several months to a year.
\section*{Results and Interpretation}
Clinical response is judged by pain reduction and improved function. Patients who do not improve may not benefit from further courses and may be candidates for other treatments.
\section*{Expected Results}
Reduced joint pain and stiffness with improved mobility and no significant adverse reaction to the injection.
\section*{Follow-up}
Patients are reassessed after the injection series, and repeated courses may be offered when symptoms recur. Failure to respond prompts consideration of corticosteroid injection, physical therapy, or joint replacement.
\section*{References}
\begin{enumerate}
  \item MedlinePlus. Viscosupplementation for the knee. U.S. National Library of Medicine; 2025.
  \item Current Medical Diagnosis and Treatment. McGraw-Hill; 2025.
\end{enumerate}

\clearpage
